\begin{figure*}[!t]
    \centering
    \includegraphics[width=1\textwidth]{fig/fig0_new.pdf}
    \caption{
    An overview of physics-informed machine learning. We review various methods of incorporating physical prior knowledge into machine learning models, ranging from strong to weak forms, such as PDEs/ODEs/SDEs, symmetry, and intuitive physics. These physical priors can be incorporated into different aspects of machine learning models, such as data, model architecture, loss function, optimizer, and inference algorithm. We also highlight different applications of physics-informed machine learning in tasks such as neural simulation, inverse problems, CV/NLP, and RL/control. Finally, we identify some significant areas for exploration in the PIML field, such as physics-informed optimizers and physics-informed inference methods.
    %Representation of physical prior from strong to weak could be divided into PDEs/ODEs/SDEs, Symmetry and intuitive physics. They could be incorporated into different parts of machine learning model like data, model architecture, loss function, optimizer and inference algorithm. Different methods are used in different types of tasks like neural simulation, inverse problems, CV/NLP and RL/control. There are still much space like physics-informed optimizer and physics-informed inference method to be explored in the field of physics-informed machine learning.
    % Representative applications in molecular dynamics, quantum chemistry, fluids dynamics and electromagnetism of physics-informed machine learning. These systems obey Schrödinger equation, Navier-Stokes equations and Maxwell's equations which are three of the most influential physical laws in science and engineering. Parts of the picture are from \cite{jumper2021highly, zhang2018deep, pathak2022fourcastnet, wu2022learning, chen2020physics}.
    }
    \vspace{-1em}
    \label{summary}
\end{figure*}
